% Part: normal-modal-logic
% Chapter: axioms-systems
% Section: normal-logics

\documentclass[../../../include/open-logic-section]{subfiles}

\begin{document}

\olfileid{nml}{prf}{nor}

\olsection{عادي مودالي منطقونه}

د مودالي !!{formula}s ځينې سټونه په اسانۍ سره د هغو
!!{formula}s په توګه نۀ شي پېژندل کېدے چې د بديهي
اصلونو له يوه سټ څخه د اشتقاق وړ دي۔ غواړو مودالي
منطقونه ښې ځانګړنې ولري۔ تر هر څه وړاندې، هر هغه څه
چې په کلاسيک بياني منطق کښې ئې !!{derive} کولے شو،
بايد لا هم !!{derivable} وي؛ البته دا په پام کښې نيسو
چې !!{formula}s اوس \iftag{prvBox}{$\Box$\iftag{prvDiamond}{
    او~}{}}{}\iftag{prvDiamond}{$\Diamond$}{} هم لرلے شي۔
د دې دپاره غواړو چې مودالي منطق د تاوتولوژيو ټول
مثالونه ولري او د مودوس پوننس لاندې تړلے وي۔

\begin{defn}
  \emph{مودالي منطق} د مودالي !!{formula}s داسې
  سټ~$\Sigma$ دے چې
  \begin{enumerate}
  \item ټولې تاوتولوژۍ لري، او
  \item د ځايناستۍ لاندې تړلے وي، يعنې که $!A \in \Sigma$
    وي، او $!D_1$، \dots، $!D_n$ !!{formula}s وي، نو
    \[
    \SSubst{!A}{\subst{!D_1}{p_1}, \dots, \subst{!D_n}{p_n}} \in \Sigma,
    \]
    \item د \emph{مودوس پوننس} لاندې تړلے وي، يعنې که $!A$
      او $!A \lif !B \in \Sigma$ وي، نو $!B \in \Sigma$ وي۔
  \end{enumerate}
\end{defn}

د مودالي منطقونو دپاره د اړيکو معناپوهنه کارولو په
موخه بايد دا هم وغواړو چې په ټولو مودالي مدلونو کښې
معتبر ټول !!{formula}s پکې شامل وي۔ څرګندېږي چې دا
غوښتنه هغه وخت پوره کېږي چې د \Ax{K}\iftag{prvDiamond}{
او~\Dual{}}{} ټول مثالونه !!{derivable} وي، او هر کله
چې !!a{formula}~$!A$ !!{derivable} وي، نو~$\Box !A$
هم وي۔ هغه مودالي منطق چې دا شرطونه پوره کړي،
\emph{عادي} بلل کېږي۔ (البته ناعادي مودالي منطقونه هم
شته، خو د اړيکو معمول مدلونه د هغو دپاره مناسب نۀ دي۔)

\begin{defn}
  مودالي منطق~$\Sigma$ \emph{عادي} دے، که دا ولري:
  \begin{align*}
    \tag{\Ax{K}} & \Box(p \lif q) \lif (\Box p \lif \Box q),
    \iftag{prvDiamond}{\\
      \tag{\Dual} & \Diamond p \liff \lnot\Box\lnot p}{}
  \end{align*}
  او د \emph{لزوم ورکولو} لاندې تړلے وي، يعنې که $!A \in
  \Sigma$ وي، نو $\Box !A \in \Sigma$ وي۔
\end{defn}

پام وکړئ چې تاوتولوژيکي استلزام د \emph{يوې نړۍ په
کښې د صدق} ساتلو دپاره په کافي اندازه دقيق دے، خو
قاعده~\Nec{} يوازې \emph{په يوه مدل کښې صدق} ساتي
(او له دې سره په يوه چوکاټ يا د چوکاټونو په يوه ټولګي
کښې اعتبار هم ساتي)۔

\begin{prop}\ollabel{prop:rk}
  هر عادي مودالي منطق د~\RK قاعدې لاندې تړلے دے:
  \begin{prooftree}
    \AxiomC{$!A_1 \lif (!A_2 \lif \cdots (!A_{n-1} \lif !A_n)\cdots)$}
    \RightLabel{\RK}
    \UnaryInfC{$\Box!A_1 \lif (\Box!A_2 \lif \cdots (\Box!A_{n-1}
      \lif \Box!A_n)\cdots).$}
  \end{prooftree}
\end{prop}

\begin{proof}
  پر~$n$ استقرا کوو: که $n = 1$ وي، نو دا قاعده
  هماغه \Nec ده، او هر عادي مودالي منطق د~\Nec لاندې
  تړلے دے۔

  اوس فرض کړئ پايله د~$n-1$ دپاره سمه ده؛ ښيو چې
  د~$n$ دپاره هم سمه ده۔

  فرض کړئ
  \begin{align*}
  & !A_1 \lif (!A_2 \lif \cdots (!A_{n-1} \lif !A_n)\cdots) \in \Sigma
  \intertext{د استقرا د فرض له مخې لرو:}
  & \Box!A_1 \lif (\Box!A_2 \lif \cdots \Box(!A_{n-1} \lif !A_n)\cdots)
  \in \Sigma
  \intertext{څرنګه چې $\Sigma$ عادي مودالي منطق دے، د~$\Ax{K}$
    ټول مثالونه لري، په ځانګړې توګه:}
  & \Box(!A_{n-1} \lif !A_n) \lif (\Box!A_{n-1} \lif \Box!A_n) \in \Sigma
  \intertext{د مودوس پوننس او د تاوتولوژۍ د مناسبو مثالونو له مخې ترلاسه کوو:}
  & \Box!A_1 \lif (\Box!A_2 \lif \cdots (\Box!A_{n-1}
  \lif \Box!A_n)\cdots) \in \Sigma. 
  \end{align*}
\end{proof}

\begin{prop}\ollabel{prop:notDiamondBot}
  هر عادي مودالي منطق~$\Sigma$، $\lnot\Diamond\lfalse$ لري۔
\end{prop}

\begin{prob}
  \olref[nml][prf][nor]{prop:notDiamondBot} ثابت کړئ۔
\end{prob}

\begin{prop}
  فرض کړئ $!A_1$، \dots، $!A_n$ !!{formula}s وي۔ بيا
  يو تر ټولو کوچنے عادي مودالي منطق~$\Sigma$ شته
  چې د~$!A_1$، \dots، $!A_n$ ټول مثالونه لري۔
  \emph{د سرچينې يادونه: په اصلي قضيه کښې «عادي» نۀ و؛
  خو ورپسې ثبوت يوازې د عادي مودالي منطقونو تقاطع
  اخلي او ورپسې تعريف هم عادي منطق ټاکي۔}
\end{prop}

\begin{proof}
  د ورکړل شويو~$!A_1$، \dots، $!A_n$ دپاره،
  $\Sigma$ د ټولو هغو عادي مودالي منطقونو تقاطع
  وټاکئ چې د~$!A_1$، \dots، $!A_n$ ټول مثالونه لري۔
  د تقاطع د سټونو کورنۍ تشه نۀ ده، ځکه چې
  $\Frm[L]$، يعنې د ټولو !!{formula}s سټ، يو داسې
  عادي مودالي منطق دے۔
\end{proof}

\begin{defn}
تر ټولو کوچني عادي مودالي منطق ته، چې $!A_1$، \dots، $!A_n$
لري، \emph{مودالي نظام} وايي، او په~$\Log{K} !A_1 \dots
!A_n$ ئې ښيو۔ تر ټولو کوچنے عادي مودالي منطق په~\Log{K}
ښودل کېږي۔
\end{defn}

\end{document}
